\section{本质的统一：一个信号，两个符号}\label{sec:unity}

\subsection{本质问题}

要判断两个策略是否“本质上相同”，需要先回答一个更基本的问题：\emph{一个交易策略究竟在赌什么}？

把任何策略的逐期损益写开，就是持仓权重与随后实现收益的内积。取期望，策略的期望收益等于
\begin{equation}
\mathbb{E}\!\left[\sum_t w_t\, r_{t+1}\right]
=\sum_t \mathrm{Cov}\!\left(w_t,\; r_{t+1}\right),
\end{equation}
其中 $r_{t+1}$ 是下一期收益，$w_t$ 是此刻的持仓（在无风险利率为零、且权重由可用信息决定的设定下，均值项可配对抵消）。这个分解给出了策略赚钱的唯一合法途径：\textbf{持仓必须与随后收益正相关}。所谓“有 alpha”，在数学上就是这个协方差为正。

于是不同策略族的区别，只能落在 $w_t$ 与 $r_{t+1}$ 各自“看什么”上：

\begin{itemize}[leftmargin=2.2em]
  \item 若 $w_t$ 取决于价格的\emph{绝对水平}或\emph{相对估值}（如市净率、展期收益率），那是价值/利差类策略——它们赌的是一阶量（水平）。
  \item 若 $w_t$ 取决于价格的\emph{历史变化}（过去一段时间的收益），那是本文讨论的对象——它们赌的是二阶量（自相关）。
\end{itemize}

这一划分是本文全部论证的起点：\textbf{趋势跟踪与均值回归同属后者，它们赌的都是符号，而不是水平。}两者都不问“这个资产现在贵不贵”，只问“它最近在往哪儿走，以及这个走势会不会持续”。

\subsection{第一性原理：唯一的自由度是符号}

设资产在 $t$ 期的收益为 $r_t$。定义\emph{信号}为过去 $L$ 期的累计收益
\begin{equation}
s_t(L)=\sum_{k=0}^{L-1} r_{t-k}.
\end{equation}
两类策略的持仓都可以写成同一个响应函数
\begin{equation}
w_t \;=\; \pm\, g\!\left(\frac{s_t(L)}{\sigma_t}\right),
\qquad g \text{ 单调不减},\ \sigma_t \text{ 为波动率尺度},
\end{equation}
其中取 $+$、$g$ 递增即趋势跟踪；取 $-$、$g$ 递减即均值回归。信号 $s_t(L)$、尺度 $\sigma_t$、窗口 $L$ 三个要素被完全共享；\textbf{唯一的自由度是那个符号}。

这个提法并不抽象，它精确对应了两类文献里的标准构造。趋势一侧，Moskowitz、Ooi \& Pedersen(2012) 的时间序列动量取
$w_t \propto \mathrm{sign}(s_t(L))/\sigma_t$，$L=12$ 个月\cite{mop2012}；截面动量则把 $s_t$ 在资产间排序后做多前列、做空后列\cite{jegadeesh1993}。反转一侧，配对交易取价差 $X_t=P^A_t-\beta P^B_t$ 对均值的\emph{负}偏离作为建仓依据，即 $w_t=-\mathrm{sign}(X_t-\mu)/\sigma$ 在阈值带外的版本\cite{gatev2006,elliott2005}；短期反转则直接用 $s_t(L)$ 的负号，$L$ 取一到四周\cite{lehmann1990,jegadeesh1990}。两组构造在符号之外逐项对应。

\subsection{同构论证：符号与回归强度是同一件事}

上面的统一仍停留在“构造形式相同”的层面。更强的同构在于：那个符号，可以严格地解释为均值回归速度的符号。

考虑对数价格 $X_t=\log P_t$ 服从 OU 过程
\begin{equation}
dX_t=\theta(\mu-X_t)\,dt+\sigma\,dW_t .
\end{equation}
$\theta>0$ 时 $X_t$ 围绕 $\mu$ 往复，偏离被拉回——这正是均值回归的连续时间载体\cite{uo1930,vasicek1977}。离散化后得到 AR(1)：
\begin{equation}
X_{t+\Delta}-\mu=e^{-\theta\Delta}\left(X_t-\mu\right)+\varepsilon_t ,
\end{equation}
自回归系数 $\phi=e^{-\theta\Delta}\in(0,1)$，自相关 $\rho_k=\phi^{k}$ 恒为正、随时间衰减，半衰期 $H=\ln 2/\theta$。

现在把 $\theta$ 的符号翻过来。若 $\theta<0$，漂移项变为 $-|\theta|(X_t-\mu)$，偏离均值不是被拉回而是被放大，价格呈指数发散——这是“趋势”的连续时间载体。此时离散系数 $\phi=e^{-|\theta|\Delta}>1$，序列非平稳，长期方差随持有期以超线性速度增长。用方差比的语言说，$\theta>0$ 对应 $VR(q)<1$，$\theta<0$ 对应 $VR(q)>1$。

\textbf{这就是本文主张的数学内核：趋势与均值回归是同一个 OU 族在两个参数半平面上的两个实例。}把 $\theta$ 从正推到负，策略从“买跌卖涨”连续地变成“追涨杀跌”。存在 $\theta=0$ 的分界，那里是随机游走，两种策略都无从获利。

这个同构带来一个立刻可用的推论：\emph{任何为均值回归设计的工具，把符号翻转后原则上都可用于趋势}。半衰期 $H$ 对应趋势的持续期 $1/|\theta|$；OU 的最优阈值带对应趋势的入场/离场带；协整关系对应“跨资产的同步趋势”。文献中“慢动量加快反转”的融合策略之所以可能，正是因为这个同构\cite{wood2021}。

\subsection{统一的统计量：方差比与自相关函数}

若符号统一了“怎么下注”，还需要一个量来刻画“事实站在哪一边”。这个量是收益的自相关结构。

设 $r_t$ 平稳，$\rho_k=\mathrm{Corr}(r_t,r_{t-k})$。$q$ 期收益的方差与 $\rho_k$ 的关系由方差比给出\cite{lomackinlay1988,campbell1997}：
\begin{equation}
VR(q)=\frac{\mathrm{Var}\!\left(\sum_{k=0}^{q-1}r_{t-k}\right)}{q\,\mathrm{Var}(r_t)}
=1+2\sum_{k=1}^{q-1}\left(1-\frac{k}{q}\right)\rho_k .
\end{equation}
$VR(q)>1$ 意味着多期收益的方差超过线性增长，即正自相关占优——这是趋势策略的生存条件；$VR(q)<1$ 意味着负自相关占优——这是反转策略的生存条件。

于是“选哪个策略”被改写为一个可以用数据回答的问题：\textbf{估出 $\rho_k$（或 $VR(q)$），看它在目标持有期上取什么符号。}

这一改写清除了很多无谓的争论。当有人问“趋势有效还是反转有效”，正确的问题是“在哪个 $q$ 上”。文献给出的答案出奇地一致——见下节。

\subsection{时间尺度不可分离：同一个自相关函数的三个区间}

把 $\rho_k$ 作为 $k$ 的函数画出来，几十年文献积累的形状是一条起伏的曲线：

\begin{itemize}[leftmargin=2.2em]
  \item \textbf{短尺度（日到周，$k\approx1$--$5$）}：$\rho_k<0$。Lehmann(1990) 的周频反转、Jegadeesh(1990) 的月频负自相关都在此区间\cite{lehmann1990,jegadeesh1990}。
  \item \textbf{中尺度（一到十二个月，$k\approx20$--$250$）}：$\rho_k>0$。Jegadeesh \& Titman(1993) 的截面动量与 Moskowitz 等(2012) 的时间序列动量在此区间\cite{jegadeesh1993,mop2012}。
  \item \textbf{长尺度（三到五年，$k\approx750$--$1250$）}：$\rho_k<0$。De Bondt \& Thaler(1985)、Fama \& French(1988)、Poterba \& Summers(1988) 的长期反转证据在此区间\cite{debondt1985,famafrench1988,poterbsummers1988}。
\end{itemize}

Cutler、Poterba \& Summers(1991) 的“投机动力学”把这种多尺度并存讲得最清楚：\emph{短期的动量证据与长期的反转证据同时成立，是因为它们瞄准了同一序列的不同频率成分}\cite{cutler1991}。这带来本文第一个可检验推论：

\begin{quote}
\itshape
若两策略共享同一信号，则在同一资产池上，趋势与反转的相对强弱应当随持有期 $L$ 单调反转——短 $L$ 上反转占优，中 $L$ 上趋势占优，长 $L$ 上又回到反转。
\end{quote}

第~\ref{sec:empirical}~节的对照回测正是为了检验这条推论（见本章末的预告）。

\subsection{统一的信息集与响应函数：为什么不能靠“减少信息”来区分}

还有一个常被忽略的统一性：两策略使用\emph{完全相同的信息集}，且都可以做到对信息量的“无损利用”。

考虑信息集 $\mathcal{F}_t$ 包含至 $t$ 为止的全部历史价格。最优的线性预测子由投影给出：$\mathbb{E}[r_{t+1}\mid\mathcal{F}_t]=\sum_k\beta_k r_{t-k+1}$。趋势策略取 $\sum_k\beta_k>0$ 的一阶近似，反转策略取 $\sum_k\beta_k<0$ 的一阶近似。二者都在用同一个 $\mathcal{F}_t$，差别只在估计出的 $\beta$ 的符号。

这一点有直接的操作含义：\textbf{不存在“因为信息不同所以结论不同”的可能性}。如果有人声称趋势策略看到的是“基本面变化”而反转策略看到的是“情绪过度”，那属于\emph{解释层}的分歧，不改变\emph{信号层}的同构；而解释层的分歧完全可以在同一个 $\rho_k$ 上并存——例如 Hong \& Stein(1999) 的信息渐进扩散模型同时给出了短期反应不足（动量）与长期过度反应（反转）\cite{hongstein1999}，Barberis 等(1998) 的投资者情绪模型也是如此\cite{bsv1998}。一个机制同时生成两个符号，这是本文在机制章要展开的内容。

\subsection{统一框架的边界}

必须诚实指出这个统一框架在哪里失效——否则它就成了过度概括。

\textbf{边界一：协整型反转不完全是时间序列反转。}配对交易的信号是\emph{跨资产的价差}偏离，而非单资产自身的历史收益。二者的可交易条件不同：协整要求两个序列存在长期均衡关系，这引入了趋势策略所没有的\emph{结构约束}。虽然形式上都可以写进响应函数，但“价差平稳”是一个比“自相关为负”更强的假设，也更易被破坏。这一差异将在第~\ref{sec:duality}~节作为第一组对偶展开。

\textbf{边界二：非线性与状态依赖。}符号统一的是\emph{一阶}描述。现实中的策略收益往往含状态依赖的非线性（趋势在低波动趋势市有效、反转在高波动区间市有效），线性自相关框架无法完全覆盖。这类现象需要 regime 切换模型处理\cite{hamilton1989,angtimmermann2012}。

\textbf{边界三：容量与成本改变了“有效”的含义。}$\rho_k$ 显著为负\emph{不等于}反转策略能赚钱。Martin \& Sch\"oneborn(2011) 的结论标题就是答案：\emph{Mean Reversion Pays, but Costs}\cite{martin2011}。统计上的正期望可能在扣除换手成本后归零。统一框架描述的是\emph{毛收益}的来源，能否落地是执行章的内容。

\textbf{边界四：波动率本身是均值回归的，且与价格方向不同构}。GARCH 类模型刻画的条件方差向长期均值回归\cite{bollerslev1986}、方差风险溢价的存在\cite{carrwu2009}与它对短期收益的预测能力\cite{bollerslev2009}，构成“波动率均值回归”这一独立家族。它的载体不是价格而是二阶矩，因此不能简单地用价格的符号翻转来刻画——这也是本文把波动率类策略列为统一框架特例而非同构对象的原因。

修正后的判断是：统一框架的适用范围是“信号层”——它精确地说明了两个策略共享什么、区别在哪；它\emph{不}声称两策略在执行、风险、容量上等价。恰恰相反，正是因为信号层如此接近，执行层的差异才成为决定成败的关键。这正是本文后半部分的重心。
